% =============================================================================
% MODEL REVISE / REWRITE — Manufacturing Enemies in Political Contests
%
% This file is a re-derivation of the *model*, *dynamic extension*, and
% *appendix proofs* that introduces a new state variable: the mobilisation
% stock M_t, accumulated from past minority counter-mobilisation, decaying
% slowly at rate rho.
%
% Architectural sources:
%   - Coate & Morris (1999, AER) — hysteresis-interval microfoundation
%     for sticky policy (their Delta vs s; here, M vs Mbar)
%   - Acharya, Lipnowski & Ramos (2025, AJPS) — promised-utility-as-state
%     Bellman with concave value function and cutoff retention rule
%   - Padro i Miquel (2007, RES) — fear-of-exclusion wedge Phi^M(M)
%     for closed-form comparative statics
%
% Notation note: M_t (italic, time-subscripted) is the stock state variable;
% subscript M on E_M, N_M continues to denote the Majority group.
%
% This file is self-contained and compiles standalone. To merge into the
% main thesis, copy Sections 2-5 to replace the current model body, copy
% Section 6 to replace the dynamic extension, and copy Section A to
% replace appendix B (proofs). The static contest of the original model
% is preserved as a special case at M=0.
% =============================================================================

\documentclass[11pt]{extarticle}
\usepackage[english]{babel}
\usepackage[breaklinks]{hyperref}
\usepackage[margin=1.2in]{geometry}
\usepackage{setspace}
\renewcommand{\baselinestretch}{1.25}
\usepackage{amsmath, amssymb, amsfonts, amsthm, mathrsfs, bm, bbm}
\usepackage{enumitem, indentfirst}
\usepackage[T1]{fontenc}
\usepackage{palatino}
\usepackage{xcolor}
\usepackage{natbib}
\usepackage{cleveref}

\newtheorem{theorem}{Theorem}
\newtheorem{lemma}{Lemma}
\newtheorem{proposition}{Proposition}
\newtheorem{corollary}{Corollary}
\newtheorem{definition}{Definition}
\newtheorem{assumption}{Assumption}
\newtheorem{remark}{Remark}
\providecommand*{\theoremautorefname}{Theorem}
\providecommand*{\lemmaautorefname}{Lemma}
\providecommand*{\propositionautorefname}{Proposition}
\providecommand*{\corollaryautorefname}{Corollary}
\providecommand*{\definitionautorefname}{Definition}
\providecommand*{\remarkautorefname}{Remark}
\providecommand*{\assumptionautorefname}{Assumption}

\hypersetup{colorlinks=true, linkcolor=blue, citecolor=blue}

\title{Manufacturing Enemies in Political Contests \\
       \large Model Revision: Mobilisation Stock and Markov Perfect Equilibrium}
\author{Haotian (George) Bai}
\date{\today}

\begin{document}
\maketitle

\begin{abstract}
\noindent This document re-derives the model, the dynamic extension, and the
appendix proofs of \emph{Manufacturing Enemies in Political Contests}, replacing
the static-contest-with-fixed-types architecture by a Markov perfect equilibrium
on a state space $(M_t, k_t)$ in which $M_t$ is a slow-decaying stock of inherited
minority mobilisation. Three technical templates are imported from the political-agency literature:
(i) the hysteresis-interval microfoundation for sticky policy from \citet{coatePolicyPersistence1999};
(ii) the promised-utility-as-state Bellman with concave value function and cutoff
retention rule from \citet{acharyaPoliticalAccountabilityMoral2025}, generalised
to a two-incumbent-type setting; and (iii) the fear-of-exclusion wedge from
\citet{padróimiquelControlPoliticiansDivided2007}, transposed to manufactured-grievance
contests as $\Phi^M(M_t)$. The static results of the original draft survive as
within-period equilibria conditional on inherited $M_{t-1}$. The dynamic extension
becomes a closed-form MPE with strategic per-period $d_t$ choice, a microfounded
sticky-environment reading, and one-line monotone comparative statics on the
conflict-trap region in $\rho$, $\lambda$, and $S_A/S_B$.
\end{abstract}

\setcounter{section}{1}

%==============================================================================
\section{Setup with Mobilisation Stock}
%==============================================================================

\subsection{Environment and the new state variable}

The economy is populated by a measure-one continuum of citizens divided into
a Majority $\mathsf M$ of share $n_M > 1/2$ and a Minority $\mathsf m$ of
share $n_m = 1 - n_M$. Two parties, the \emph{Hawk} ($A$) and the \emph{Dove}
($B$), compete for the support of the Majority electorate. Parties are
distinguished by a two-dimensional technological endowment $(S_k, \bar g_k)$ with
$S_A > S_B$ and $\bar g_A < \bar g_B$. As in the original draft, the
public-good policy binds at the technological ceiling $g_k^* = \bar g_k$
in every period, so only the symbolic-policy choice $d_k \ge 0$ carries
strategic content.

\paragraph{The mobilisation stock $M_t$.}
The new state variable is a scalar $M_t \in \mathbb{R}_+$ — the inherited
\emph{mobilisation stock} of the Minority. It captures accumulated grievance
and organisational commitment that persist beyond a single period: if minority
counter-mobilisation has been intense in recent periods, $M_t$ is high; if
not, $M_t$ is near zero.

The stock evolves according to
\begin{equation} \label{eq:M_evolution}
    M_{t+1} \;=\; (1-\rho)\,M_t \;+\; \rho\,E_m^*(d_t, S_{k_t}; M_t),
\end{equation}
where $\rho \in (0, 1)$ is the depreciation rate, $E_m^*(d_t, S_{k_t}; M_t)$ is
the within-period equilibrium minority effort (defined in \autoref{lem:contest_M}
below), and $k_t \in \{A, B\}$ is the period-$t$ incumbent. The initial stock
$M_0 \ge 0$ is exogenous. We refer to $1 - \rho$ as the \emph{stock persistence}.

\subsection{Within-period contest with stock-augmented stake}

The Majority's within-period collective-action problem is unchanged:
each citizen $i \in \{1, \dots, N_M\}$ chooses enforcement effort $e_i \ge 0$
at cost $\frac{1}{2}e_i^2$, deriving public benefit
$\Lambda \cdot \Pi_k(E_M, E_m)$ from the Tullock contest success function
\begin{equation} \label{eq:CSF_M}
    \Pi_k(E_M, E_m) \;=\; \frac{S_k\,E_M}{S_k\,E_M + E_m}.
\end{equation}

The Minority's collective-action problem is modified to absorb the inherited
stock as a \emph{stake amplifier}. Citizen $j \in \{1, \dots, N_m\}$ chooses
effort $e_j \ge 0$ at private cost $\frac{1}{2\kappa}e_j^2$ to maximise
\begin{equation} \label{eq:minority_M}
    V_j \;=\; \phi\,\big(d^\sigma + \xi M\big)\,\bigl(1 - \Pi_k(E_M, E_m)\bigr)
    \;-\;\frac{1}{2\kappa} e_j^2,
\end{equation}
where $\xi > 0$ is the \emph{stock-stake elasticity} and $\sigma > 0$ the
grievance-elasticity-of-policy parameter. The new term $\xi M$ captures inherited
grievance: even at $d = 0$, the minority has a non-zero stake whenever
$M > 0$, because past mobilisation has built organisational capacity and
identity salience that current peace cannot immediately undo.

\begin{remark} \label{rem:stake_amplifier}
The stock enters the stake additively: the \emph{effective} provocation
intensity is $\hat d^\sigma := d^\sigma + \xi M$. The original model is the
special case $\xi = 0$, equivalently $M = 0$ in every period.
This formulation has three consequences:
(i) the within-period contest collapses to the original model's algebra
under the substitution $d^\sigma \to \hat d^\sigma$, preserving all closed forms;
(ii) at $d_t = 0$ with $M_t > 0$, the contest is \emph{not} dormant — the
minority mobilises in proportion to inherited grievance;
(iii) the Hawk's optimal provocation $d_A^*$ becomes a \emph{decreasing}
function of $M$ — more inherited grievance substitutes for current
provocation, a feature that becomes load-bearing in the dynamic equilibrium.
\end{remark}

\subsection{Joint contest equilibrium with stock}

Let $\alpha := \sqrt{\Lambda N_M / (\kappa\phi N_m)}$ and
$\beta := (\kappa \phi N_m \Lambda N_M)^{1/4}$ exactly as in the original
draft. The within-period contest with inherited stock has the following
characterisation.

\begin{lemma}[\textbf{Joint contest equilibrium with stock}] \label{lem:contest_M}
Fix $(d, S_k, M) \in [0,\infty) \times (0, \infty) \times [0, \infty)$ and write
\begin{equation} \label{eq:dhat}
    \hat d^\sigma \;:=\; d^\sigma + \xi M.
\end{equation}
If $\hat d^\sigma = 0$ (equivalently $d = 0$ and $M = 0$), the contest is
dormant: $E_M^* = E_m^* = 0$ and $\lim_{\hat d \downarrow 0} \Pi_k^* = 1$.
If $\hat d^\sigma > 0$, the unique symmetric Nash equilibrium is
\begin{equation} \label{eq:EmEM_M}
    E_m^*(d, S_k; M) \;=\; \frac{\beta^2\, S_k^{1/2}\,\hat d^{\,3\sigma/4}}{S_k\alpha + \hat d^{\,\sigma/2}},
    \qquad
    E_M^*(d, S_k; M) \;=\; \alpha\,\beta^2\,\frac{S_k^{1/2}\,\hat d^{\,\sigma/4}}{S_k\alpha + \hat d^{\,\sigma/2}},
\end{equation}
and the equilibrium security probability is
\begin{equation} \label{eq:Pi_M}
    \Pi_k^*(d, S_k; M) \;=\; \frac{S_k\alpha}{S_k\alpha + \hat d^{\,\sigma/2}}.
\end{equation}
The function $\Pi_k^*$ is strictly increasing in $S_k$, strictly decreasing
in $\hat d$ (hence in both $d$ and $M$), and the limits
$\lim_{\hat d \downarrow 0}\Pi_k^* = 1$, $\lim_{\hat d \uparrow \infty}\Pi_k^* = 0$ hold.
\end{lemma}

\begin{proof}[Proof sketch] Apply the substitution $d^\sigma \to \hat d^\sigma = d^\sigma + \xi M$ in the FOC system of \autoref{eq:minority_M} and the unchanged Majority FOC. The algebra of the original Joint Contest Equilibrium Lemma carries through unchanged. Full proof in \autoref{app:proof:contest_M}.
\end{proof}

The Security Premium is now a function of both $d$ and the inherited stock:
\begin{equation} \label{eq:DPi_M}
    \Delta\Pi(d; M) \;:=\; \Pi_A^*(d, S_A; M) - \Pi_B^*(d, S_B; M)
    \;=\; \frac{S_A\alpha}{S_A\alpha + \hat d^{\,\sigma/2}} - \frac{S_B\alpha}{S_B\alpha + \hat d^{\,\sigma/2}}.
\end{equation}

\begin{lemma}[\textbf{Inverted-U Security Premium with stock and migrating peak}] \label{lem:invU_M}
Suppose $S_A > S_B > 0$. The Security Premium $\Delta\Pi(\cdot; M)$ as a
function of current $d$ is single-peaked on $[0, \infty)$ for every fixed
$M \ge 0$. The unique maximiser is
\begin{equation} \label{eq:dpeak_M}
    d^{\mathrm{peak}}(M) \;=\;
    \begin{cases}
        \big(\alpha^2 S_A S_B - \xi M\big)^{1/\sigma} & \text{if } M < \alpha^2 S_A S_B / \xi, \\[0.4em]
        0 & \text{if } M \ge \alpha^2 S_A S_B / \xi.
    \end{cases}
\end{equation}
The peak value is unchanged from the original model:
\begin{equation} \label{eq:DPi_peak_M}
    \max_{d \ge 0}\Delta\Pi(d; M) \;=\; \frac{\sqrt{S_A} - \sqrt{S_B}}{\sqrt{S_A} + \sqrt{S_B}},
\end{equation}
attained at $d^{\mathrm{peak}}(M)$ for any $M < \alpha^2 S_A S_B / \xi$, and
attained at the boundary $d = 0$ for higher $M$.
\end{lemma}

\autoref{lem:invU_M} contains the central new comparative static of the
revised model: \emph{the optimal provocation peak migrates leftward as the
inherited stock grows}. At $M = 0$ the peak is at the original
$d^{\mathrm{peak}}(0) = (\alpha\sqrt{S_A S_B})^{2/\sigma}$; at sufficiently large
$M$, the peak coincides with $d = 0$ — the inherited grievance has already
done the contest-activation work, and the Hawk's marginal incentive to
provoke vanishes. This is the structural source of the \emph{paradox of
restraint} in the dynamic equilibrium: a sufficiently entrenched Hawk
chooses minimal current provocation precisely because the stock is doing
the work for him.

%==============================================================================
\section{Voter Preferences and Within-Period Equilibrium}
%==============================================================================

\subsection{Per-period payoffs}

The Majority representative voter's per-period utility under incumbent
$k \in \{A, B\}$ at inherited stock $M$ and current policy $d$ is
\begin{equation} \label{eq:utility_M}
    U^k(d; M) \;=\; \lambda\, \Pi_k^*(d, S_k; M) \;+\; (1-\lambda)\, u(\bar g_k)
    \;-\; \chi_{\mathrm{econ}}(d) \;-\; \chi_{\mathrm{mod}}(d)\,\mathbb{1}[k = A]\,\mathbb{1}[d > 0],
\end{equation}
where $\lambda \in (0, 1)$ is the (exogenous) weight on security, $u(\cdot)$
is strictly increasing in the public good, $\chi_{\mathrm{econ}}, \chi_{\mathrm{mod}}$
are convex, increasing, with $\chi_{\mathrm{econ}}(0) = \chi_{\mathrm{mod}}(0) = 0$
(unchanged from original). The economic deficit is
$\Delta u := u(\bar g_B) - u(\bar g_A) > 0$ as before.

\subsection{Within-period election under sticky environment}

A central simplification of the original draft was the ``sticky-environment''
reading: voters compare the two parties at the \emph{incumbent's} policy $d$,
not at counterfactual policies. With $M_t$ as a state variable, this is
\emph{microfounded} rather than assumed. Within period $t$, the inherited
stock $M_{t-1}$ has already determined the contest's stake amplifier
$\xi M_{t-1}$, which both parties face. The current $d_t$ is the only
strategic margin; voters compare $U^A(d_t; M_{t-1})$ and $U^B(d_t; M_{t-1})$.

Adding an idiosyncratic taste shock $\xi_i \sim \Phi$ in favour of the Hawk,
voter $i$ supports $A$ iff
$U^A(d_t; M_{t-1}) - U^B(d_t; M_{t-1}) + \xi_i \ge 0$. The swing-voter
indifference condition is
\begin{equation} \label{eq:swing_M}
    \hat\xi(d, M) \;=\; (1-\lambda)\,\Delta u \;+\; \chi_{\mathrm{mod}}(d)\,\mathbb{1}[d > 0]
    \;-\; \lambda\,\Delta\Pi(d; M).
\end{equation}
Party $A$ wins iff $\Phi(-\hat\xi(d, M)) > 1/2$.

\begin{remark}[Microfoundation of sticky policy] \label{rem:microfound_sticky}
At $d = 0$ and $M = 0$, both parties give $\Pi^* = 1$, so $\Delta\Pi = 0$ and
the security dimension is electorally invisible — the voter elects on the
economic dimension and the Dove wins (whenever $\Delta u > 0$). At $d = 0$
and $M > 0$, $\Delta\Pi(0; M) > 0$: the Hawk's $S_A$-advantage is
electorally salient \emph{even with no current provocation}, purely
because of inherited grievance. This is the structural microfoundation
the original draft relied on assuming. Following \citet{coatePolicyPersistence1999},
the inherited stock plays the role of irreversible private investment that
locks in the policy regime.
\end{remark}

\subsection{Hawk's static (within-period) provocation problem}

Given inherited $M$, the Hawk-incumbent's optimal $d$ solves
\begin{equation} \label{eq:Hawk_FOC_M}
    d_A^*(M) \;\in\; \arg\max_{d \ge 0}\;\big[\lambda\,\Delta\Pi(d; M) - \chi_{\mathrm{mod}}(d)\big].
\end{equation}

Define the within-period \emph{net electoral return} to provocation as
$F(d; M) := \lambda\,\Delta\Pi(d; M) - \chi_{\mathrm{mod}}(d)$, and the
within-period provocation cutoff as
\begin{equation} \label{eq:du_bar_M}
    \overline{\Delta u}(M) \;:=\; \frac{1}{1-\lambda}\,\max_{d \ge 0} F(d; M).
\end{equation}

\begin{proposition}[\textbf{Within-period strategic provocation with stock}] \label{prop:provocation_M}
Under the conditions on $\chi_{\mathrm{mod}}$ in the original draft (continuously
differentiable, strictly increasing, strictly convex, $\chi_{\mathrm{mod}}(0) = 0$,
$\lim_{d\uparrow\infty}\chi_{\mathrm{mod}}(d) = \infty$):
\begin{enumerate}[label=(\roman*)]
    \item For each $M \ge 0$, $F(\cdot; M)$ attains a finite maximum at some
    $d_A^*(M) \ge 0$, and $\overline{\Delta u}(M) \ge 0$.
    \item A Hawk victory is feasible iff $\Delta u \le \overline{\Delta u}(M)$.
    \item The cutoff is non-decreasing in $M$:
    $\partial \overline{\Delta u}(M)/\partial M \ge 0$, with strict inequality
    whenever $d_A^*(M) > 0$ is interior.
    \item The optimal provocation $d_A^*(M)$ is non-increasing in $M$:
    when inherited grievance grows, the Hawk needs less current provocation
    to maintain its electoral advantage.
\end{enumerate}
\end{proposition}

Part (iii) is the analogue of Coate--Morris's central comparative static —
the willingness to pay for the policy is increasing in inherited stock —
transposed to the manufactured-enemy setting. Part (iv) is the
\emph{paradox of restraint}: a Hawk in a high-$M$ steady state may quite
literally choose $d_t = 0$ in equilibrium and still hold the security
premium intact. The original draft could not deliver this prediction.

%==============================================================================
\section{Dynamic Extension: Markov Perfect Equilibrium}
%==============================================================================

\subsection{State, timing, and strategies}

The political game is dynamic with discrete time $t = 0, 1, 2, \dots$ and an
infinite horizon. The state at the start of period $t$ is
$(M_t, k_t) \in \mathbb{R}_+ \times \{A, B\}$, where $k_t$ is the incumbent
inherited from period $t-1$'s election outcome and $M_t$ is the inherited
mobilisation stock. (We write $M_t$ for the state at the \emph{start} of
period $t$; the original draft's notation has $M_{t-1}$ in the same role.)

Within each period:
\begin{enumerate}[label=\textbf{Stage \arabic*:}, leftmargin=*, itemsep=0.2em]
    \item \textbf{Policy.} Incumbent $k_t$ chooses $d_t \ge 0$.
    \item \textbf{Contest.} The within-period contest equilibrium
    $(E_M^*, E_m^*, \Pi_k^*)$ realises at $(d_t, S_{k_t}, M_t)$ per \autoref{lem:contest_M}.
    \item \textbf{Election.} The Majority electorate observes a noisy signal
    of period-$t$ outcomes and chooses $k_{t+1} \in \{A, B\}$.
    \item \textbf{Stock update.} $M_{t+1} = (1-\rho)M_t + \rho E_m^*(d_t, S_{k_t}; M_t)$.
\end{enumerate}

A \emph{Markov strategy} for the incumbent is a measurable function
$d_k : \mathbb{R}_+ \to \mathbb{R}_+$ specifying provocation at each $M$ if
party $k$ holds office. A \emph{Markov retention rule} for the voter is a
measurable function $\rho^v : \mathbb{R}_+ \times \{A, B\} \to [0, 1]$
specifying the probability of retaining the incumbent at each state.

\subsection{Voter Bellman and the cutoff retention rule}

We follow \citet{acharyaPoliticalAccountabilityMoral2025} in solving the
voter's problem recursively. Let $V^v(M, k)$ denote the voter's expected
discounted utility at state $(M, k)$. Given the incumbent's strategy
$d_k(\cdot)$ and the law of motion \autoref{eq:M_evolution}:
\begin{equation} \label{eq:voter_bellman}
    V^v(M, k) \;=\; U^k(d_k(M); M) \;+\; \delta\,\big[\rho^v(M, k)\,V^v(M', k) + (1 - \rho^v(M, k))\,V^v(M', \tilde k)\big],
\end{equation}
where $M' = (1-\rho)M + \rho E_m^*(d_k(M), S_k; M)$, $\tilde k = B$ if $k = A$
and vice versa, and $\delta \in (0, 1)$ is the voter's discount factor.

Optimal retention compares the two regime-specific value functions:
$\rho^v(M, k) = 1$ if $V^v(M, k) > V^v(M, \tilde k)$ and $0$ otherwise.
Define the \emph{voter wedge}
\begin{equation} \label{eq:wedge}
    \Phi^M(M) \;:=\; V^v(M, A) - V^v(M, B).
\end{equation}
The voter retains the Hawk at state $(M, A)$ iff $\Phi^M(M) \ge 0$. The wedge
is the structural transposition of \citet{padróimiquelControlPoliticiansDivided2007}'s
fear-of-exclusion wedge to the manufactured-grievance setting: it bundles
discounting (through $V^v$), incumbency stickiness (through retention
probability), and cross-regime utility gap (through $V^v(M, A) - V^v(M, B)$)
into a single scalar.

\begin{lemma}[\textbf{Concavity and monotonicity of $V^v$}] \label{lem:Vv_concave}
Under the conditions of \autoref{prop:provocation_M}:
\begin{enumerate}[label=(\roman*)]
    \item $V^v(\cdot, A)$ and $V^v(\cdot, B)$ are continuous and uniformly
    bounded on $[0, \infty)$.
    \item $V^v(M, B)$ is non-increasing in $M$: under a Dove incumbent at
    $d_B = 0$, higher inherited grievance is unambiguously bad for the
    voter's per-period payoff, and improvement requires waiting for
    decay.
    \item $V^v(M, A)$ is single-peaked in $M$: at low $M$ the Hawk's
    advantage is electorally invisible (low security premium); at high
    $M$ the cumulative cost $\chi_{\mathrm{econ}}(d_A^*(M))$ + the realised
    security loss outpaces the security-premium gain.
    \item Consequently, the voter wedge $\Phi^M(M) = V^v(M, A) - V^v(M, B)$
    is single-crossing: there exists a unique $\bar M \in (0, \infty]$ such
    that $\Phi^M(M) < 0$ for $M < \bar M$ and $\Phi^M(M) \ge 0$ for $M \ge \bar M$.
\end{enumerate}
\end{lemma}

The single-crossing wedge yields the cutoff retention rule. The Acharya
et al. template for our setting:

\begin{proposition}[\textbf{Cutoff retention rule}] \label{prop:cutoff}
A Markov retention rule $\rho^v$ is part of a Markov Perfect Equilibrium
if and only if
\begin{equation} \label{eq:cutoff}
    \rho^v(M, A) \;=\;
    \begin{cases}
        1 & \text{if } M \ge \bar M, \\
        0 & \text{if } M < \bar M,
    \end{cases}
    \qquad
    \rho^v(M, B) \;=\;
    \begin{cases}
        1 & \text{if } M \le \underline M, \\
        0 & \text{if } M > \underline M,
    \end{cases}
\end{equation}
where $\bar M = \inf\{M : \Phi^M(M) \ge 0\}$ and $\underline M = \sup\{M : \Phi^M(M) \le 0\}$.
The two thresholds satisfy $\underline M \le \bar M$, with strict inequality
generic in $\rho$.
\end{proposition}

The thresholds $\bar M$ and $\underline M$ play the same role as Coate--Morris's
$\Delta$ vs $s$: the interval $[\underline M, \bar M]$ is the
\emph{hysteresis region}, where retention depends on which incumbent is
currently in office rather than on the state $M$ alone.

\subsection{Existence of MPE}

\begin{proposition}[\textbf{Existence of pure-strategy MPE}] \label{prop:MPE_exist}
Under the regularity conditions on $\chi_{\mathrm{econ}}$, $\chi_{\mathrm{mod}}$,
and the contest primitives, a pure-strategy Markov Perfect Equilibrium
$(d_A^\star, d_B^\star, \rho^v)$ exists in which:
\begin{enumerate}[label=(\roman*)]
    \item $d_B^\star(M) = 0$ for all $M \ge 0$.
    \item $d_A^\star(M)$ is the within-period optimum
    $\arg\max_d F(d; M)$ from \autoref{prop:provocation_M}, weakly decreasing in $M$.
    \item $\rho^v$ is the cutoff rule of \autoref{prop:cutoff}.
\end{enumerate}
The voter value functions $V^v(\cdot, A), V^v(\cdot, B)$ are the unique
fixed point of the Bellman operator on the Banach space of bounded
continuous functions on $[0, \infty)$.
\end{proposition}

The Dove's corner $d_B^\star = 0$ now has a different interpretation than in
the original draft: even though the inherited stock $M$ keeps the Minority
mobilising and pushes $\Pi_B^* < 1$, the Dove still chooses $d = 0$ because
provoking would (i) push $\Pi_B^*$ even lower and (ii) generate
$\chi_{\mathrm{mod}}$ on the Dove. The Dove cannot benefit from raising the
security dimension because $S_B < S_A$. The corner survives the introduction
of $M$.

%==============================================================================
\section{Long-Run Outcomes: Conflict Trap, Cycles, and the Paradox}
%==============================================================================

\subsection{Conflict trap restated}

\begin{proposition}[\textbf{Conflict trap with mobilisation stock}] \label{prop:trap_M}
Suppose the model parameters $(\rho, \delta, \lambda, S_A, S_B, \xi, \kappa, \phi, \sigma, \Lambda, N_M, N_m)$
are such that $\bar M < \infty$ and $d_A^\star(\bar M) > 0$. Then there
exists $\bar\rho \in (0, 1)$ such that for all $\rho < \bar\rho$, the set
\begin{equation} \label{eq:trap_set}
    \mathcal{T} \;:=\; \big\{(M, k) : k = A,\ M \ge \bar M\big\}
\end{equation}
is forward-invariant under the equilibrium dynamics: starting from any
$(M_0, A) \in \mathcal{T}$, $(M_t, k_t) \in \mathcal{T}$ for all $t \ge 0$
almost surely. Within the trap, the equilibrium policy converges to a
stationary level $d_A^\infty > 0$ (or $d_A^\infty = 0$ in the
\emph{paradox-of-restraint} regime $M^\infty \ge \alpha^2 S_A S_B / \xi$)
and the Hawk is retained in every period.
\end{proposition}

\autoref{prop:trap_M} differs from the v1 conflict trap (which required
forward-looking voters and an exogenous threat shock) in two ways:
\begin{enumerate}[label=(\roman*)]
    \item No exogenous threat shock is needed. The trap is generated
    endogenously by the strategic interaction between Hawk's
    provocation and slow stock decay. Voter myopia ($\delta = 0$) is
    sufficient.
    \item ``Permanent capacity loss'' from a Dove victory is no longer
    counterfactual but \emph{quantitative}: switching to Dove leaves
    inherited $M_t$ unchanged in the short run, but Dove's $d_B = 0$
    causes $M$ to decay at rate $\rho$. Voter discounting makes this
    transition welfare-decreasing if $\rho$ is small relative to the
    Hawk's per-period security advantage. Below the cutoff $\bar M$, the
    voter prefers Dove; above, Hawk.
\end{enumerate}

The paradox of restraint within the trap is the most distinctive new
prediction: a sufficiently entrenched Hawk regime can equilibrate at $d_t = 0$
indefinitely, with the contest sustained entirely by inherited grievance.
This rationalises observed regimes where the security-coded party governs
without enacting visibly new symbolic policy — Iran's hijab regime in its
mature phase, India's BJP after 2019 — without abandoning the manufactured-enemies
mechanism.

\subsection{Political cycles restated}

When $\rho$ is large (fast stock decay), the trap region shrinks. Under
myopic voting and an exogenous threat process $\omega_t \in [0, \bar\omega]$
that adds to the within-period stake, the equilibrium exhibits regime
oscillations.

\begin{proposition}[\textbf{Political cycles with mobilisation stock}] \label{prop:cycles_M}
Suppose voters are myopic ($\delta \to 0$) and the per-period stake includes
an exogenous threat $\omega_t$ with full support on $[0, \bar\omega]$.
Suppose further that
$\bar\rho < \rho < 1$ (fast decay regime). Then:
\begin{enumerate}[label=(\roman*)]
    \item Conditional on a Hawk incumbent, the probability of a Hawk-to-Dove
    transition is strictly decreasing in current $\omega_t$ and in inherited
    $M_t$.
    \item Conditional on a Dove incumbent, the probability of a Dove-to-Hawk
    transition is strictly increasing in $\omega_t$, in $\xi$, and in the
    minority's structural mobilisation potential $\kappa\phi$.
    \item Both transitions occur with positive probability along the
    realised path, yielding recurrent regime switching.
\end{enumerate}
\end{proposition}

The substantive distinction from the v1 cycles result is that $d_t$ is now
\emph{endogenously oscillating} via $d_A^\star(M_t)$, rather than mechanically
flipping between fixed type values. Under a Hawk-to-Dove transition, the
inherited stock $M_t$ decays gradually under $d_B = 0$; the contest
softens. Under a Dove-to-Hawk transition, the new Hawk inherits a low
$M_t$ and chooses a high $d_t$ to re-activate the contest; the cycle
restarts.

\subsection{Paradox of competence}

\begin{corollary}[\textbf{Paradox of competence with stock}] \label{cor:paradox_M}
Under myopic voting, a first-order stochastic decrease in observed disorder
$\tilde v_t = (1 - \Pi_{k_t}^*) + \varepsilon_t^v$ under a Hawk incumbent
weakly decreases the perceived security premium and weakly increases the
hazard of a Hawk-to-Dove transition.
\end{corollary}

The corollary carries through unchanged from the v1 paradox of competence:
quieter streets shrink the perceived security premium and raise the Hawk's
defeat hazard. The new dynamic mechanism does not affect the underlying
inference problem.

%==============================================================================
\section{Comparative Statics on the Trap Region}
%==============================================================================

The fear-wedge $\Phi^M$ delivers one-line comparative statics on the
conflict-trap region.

\begin{proposition}[\textbf{Comparative statics on $\bar M$}] \label{prop:cs_trap}
The conflict-trap threshold $\bar M$ satisfies
\begin{equation} \label{eq:cs_signs}
    \frac{\partial \bar M}{\partial \rho} \;>\; 0,
    \qquad \frac{\partial \bar M}{\partial (S_A/S_B)} \;<\; 0,
    \qquad \frac{\partial \bar M}{\partial \lambda} \;<\; 0,
    \qquad \frac{\partial \bar M}{\partial \Delta u} \;>\; 0.
\end{equation}
That is, the trap region (states with $M \ge \bar M$ that are forward-invariant) shrinks in faster decay $\rho$ and in larger Dove economic advantage $\Delta u$, and expands in larger Hawk security advantage $S_A/S_B$ and in larger weight $\lambda$ on security.
\end{proposition}

These four monotone comparative statics replace the original draft's
implicit reliance on numerical examples. Each runs through the wedge
$\Phi^M$ and the implicit-function theorem applied to
$\Phi^M(\bar M; \rho, \lambda, S_A/S_B, \Delta u) = 0$.

%==============================================================================
\section{Discussion: How This Differs from the Original Draft}
%==============================================================================

The original draft has three load-bearing assumptions that this revision
either eliminates or microfounds:

\begin{enumerate}
    \item \textbf{Sticky-environment swing-voter equation.} The original
    draft assumes both parties are evaluated at the incumbent's $d$ — see
    the v1 footnote at the swing-voter equation. This is now microfounded:
    inherited $M_{t-1}$ has already determined the within-period stake
    amplifier $\xi M_{t-1}$, which is exogenous to the current incumbent's
    choice.

    \item \textbf{Fixed party types $d_A > d_B$ in dynamic.} The v1 dynamic
    extension assumed $d_A > d_B$ as exogenous types, contradicting the
    static result $d_B^* = 0$. The revised dynamic extension has both
    parties making strategic per-period choices: $d_B^\star(M) = 0$
    (consistent with static), $d_A^\star(M)$ a continuous-valued strategic
    function (resolving the contradiction).

    \item \textbf{Permanent security-capacity loss from regime change.}
    The v1 trap required voters to expect future exogenous threats that
    Dove's $S_B$ cannot manage. The revised trap uses inherited $M$
    directly: under Dove, $d_B = 0$ and $M$ decays at rate $\rho$, but in
    the meantime the Hawk's $S_A$ would have been valuable. Voter
    discounting + slow decay = trap, with no exogenous threat process
    needed.
\end{enumerate}

In addition, the revised draft contributes a named state-variable type
to the political-agency literature: the \emph{adversary's mobilisation
stock}. Existing political-agency models work with two types of state
variables: voter beliefs about the politician's type or future actions
(e.g., \citealp{acharyaPoliticalAccountabilityMoral2025};
\citealp{banksRepeatedElectionsModel1993}), and the politician's own
reputation or commitment stock (e.g., \citealp{coatePolicyPersistence1999};
\citealp{padróimiquelControlPoliticiansDivided2007}). The mobilisation
stock $M_t$ is a third type: it accumulates from the adversary's response
to the agent's strategic action, and yet the strategic agent has an
interest in deliberately accumulating it because doing so activates the
very contest dimension on which his comparative advantage operates. This
is the formal version of ``manufactured persistence.''

%==============================================================================
\appendix
\section{Proofs} \label{app:proofs}
%==============================================================================

\subsection{Proof of \autoref{lem:contest_M}} \label{app:proof:contest_M}

\begin{proof}
The Majority FOC of the original draft is unchanged:
\[
    e_i \;=\; \Lambda \,\frac{S_k E_m}{(S_k E_M + E_m)^2}.
\]
The Minority FOC under \autoref{eq:minority_M} replaces the original's
$\phi d^\sigma$ with $\phi(d^\sigma + \xi M)$:
\[
    e_j \;=\; \kappa\,\phi(d^\sigma + \xi M)\,\frac{S_k E_M}{(S_k E_M + E_m)^2}.
\]
Define $\hat d^\sigma := d^\sigma + \xi M$. The system is now identical to the
original draft's contest equilibrium with $d^\sigma$ replaced by
$\hat d^\sigma$ throughout. Dividing the symmetric FOCs $E_M/N_M = e_i$
and $E_m/N_m = e_j$ yields
\[
    \frac{E_M/N_M}{E_m/N_m} \;=\; \frac{\Lambda E_m}{\kappa\phi\hat d^\sigma E_M},
\]
hence $\kappa\phi\hat d^\sigma N_m E_M^2 = \Lambda N_M E_m^2$, which gives
$E_M / E_m = \alpha\hat d^{-\sigma/2}$. Substituting into the Majority FOC
yields $E_m^*$ as in \autoref{eq:EmEM_M}. The Tullock CSF then gives
$\Pi_k^* = S_k E_M^* / (S_k E_M^* + E_m^*) = S_k\alpha / (S_k\alpha + \hat d^{\sigma/2})$
as in \autoref{eq:Pi_M}. Comparative statics in $S_k$, $d$, and $M$ follow
from differentiating these closed forms. The dormant case $\hat d = 0$
follows from continuity: $\lim_{\hat d \downarrow 0} E_m^* = 0$,
$\lim_{\hat d \downarrow 0} E_M^* = 0$,
$\lim_{\hat d \downarrow 0}\Pi_k^* = 1$.
\end{proof}

\subsection{Proof of \autoref{lem:invU_M}} \label{app:proof:invU_M}

\begin{proof}
By \autoref{eq:DPi_M}, write $\Delta\Pi(d; M) = h(\hat d^{\sigma/2})$ where
\[
    h(x) := \frac{S_A\alpha}{S_A\alpha + x} - \frac{S_B\alpha}{S_B\alpha + x}.
\]
The function $h$ is the same as in the original draft's Inverted-U Security
Premium lemma. It is strictly positive on $(0, \infty)$
when $S_A > S_B$, vanishes at $0$ and $\infty$, and is single-peaked at
$x^{\mathrm{peak}} = \alpha\sqrt{S_A S_B}$ with peak value
$h(x^{\mathrm{peak}}) = (\sqrt{S_A} - \sqrt{S_B})/(\sqrt{S_A} + \sqrt{S_B})$.

Translating back: $\Delta\Pi(d; M)$ is single-peaked in $\hat d^{\sigma/2}$
at $\alpha\sqrt{S_A S_B}$, equivalently $\hat d^\sigma = \alpha^2 S_A S_B$,
equivalently $d^\sigma + \xi M = \alpha^2 S_A S_B$. Solving for $d$:
\begin{enumerate}
    \item If $\xi M < \alpha^2 S_A S_B$, the peak is interior at
    $d^{\mathrm{peak}}(M) = (\alpha^2 S_A S_B - \xi M)^{1/\sigma}$.
    \item If $\xi M \ge \alpha^2 S_A S_B$, the peak is at the boundary $d = 0$
    because $\hat d^\sigma$ already exceeds $\alpha^2 S_A S_B$ from the
    stock alone, and any further $d > 0$ moves $\hat d^{\sigma/2}$ past
    $\alpha\sqrt{S_A S_B}$ into the decreasing part of $h$.
\end{enumerate}
The peak value is $h(\alpha\sqrt{S_A S_B})$ in case (1) and
$h(\sqrt{\xi M}) < h(\alpha\sqrt{S_A S_B})$ in case (2). Within case (1),
the peak value is $M$-invariant; within case (2), it strictly decreases in $M$.
\end{proof}

\subsection{Proof of \autoref{prop:provocation_M}} \label{app:proof:provocation_M}

\begin{proof}
\textit{(i)} $F(d; M) = \lambda\Delta\Pi(d; M) - \chi_{\mathrm{mod}}(d)$ is
continuous on $[0, \infty)$ for each $M$, with $F(0; M) = \lambda\Delta\Pi(0; M)$
finite and $\lim_{d\uparrow\infty}F(d; M) = -\infty$ since $\chi_{\mathrm{mod}}(d) \to \infty$
and $\Delta\Pi$ is bounded by 1. Hence the maximum is attained at some
finite $d_A^*(M)$. $\overline{\Delta u}(M) = F(d_A^*(M); M)/(1-\lambda) \ge F(0; M)/(1-\lambda) \ge 0$
because $\Delta\Pi \ge 0$ everywhere.

\textit{(ii)} The Hawk wins iff the swing-voter threshold
$\hat\xi(d_A^*(M); M) \le 0$, i.e., $F(d_A^*(M); M) \ge (1-\lambda)\Delta u$,
i.e., $\overline{\Delta u}(M) \ge \Delta u$.

\textit{(iii)} By the envelope theorem,
\[
    \frac{\partial}{\partial M}\overline{\Delta u}(M)
    \;=\; \frac{1}{1-\lambda}\,\frac{\partial F(d_A^*(M); M)}{\partial M}
    \;=\; \frac{\lambda}{1-\lambda}\,\frac{\partial \Delta\Pi(d_A^*(M); M)}{\partial M}.
\]
The partial derivative of $\Delta\Pi$ with respect to $M$ at any
$d \in (0, d^{\mathrm{peak}}(0))$ is positive on the increasing branch of
the inverted-U: increasing $M$ moves $\hat d^{\sigma/2}$ closer to the peak.
On the decreasing branch, the sign flips and $\partial\Delta\Pi/\partial M < 0$.
For interior $d_A^*(M) \in (0, d^{\mathrm{peak}}(M))$, the optimum lies on
the increasing branch (where the marginal benefit equals
$\chi_{\mathrm{mod}}'$); hence
$\partial\overline{\Delta u}/\partial M > 0$.

\textit{(iv)} Differentiate the Hawk's FOC
$\lambda\,\partial\Delta\Pi/\partial d|_{d_A^*(M); M} = \chi_{\mathrm{mod}}'(d_A^*(M))$
implicitly with respect to $M$. By the convexity of $\chi_{\mathrm{mod}}$
and the supermodularity of $\Delta\Pi$ in $(d, -M)$ on the relevant region:
$\partial d_A^*/\partial M < 0$. Details: differentiate
\[
    \lambda\,\frac{\partial \Delta\Pi(d; M)}{\partial d}\bigg|_{d=d_A^*(M)} \;=\; \chi_{\mathrm{mod}}'(d_A^*(M))
\]
with respect to $M$ to get
\[
    \lambda\,\frac{\partial^2 \Delta\Pi}{\partial d^2}\,\frac{\partial d_A^*}{\partial M} + \lambda\,\frac{\partial^2 \Delta\Pi}{\partial d\,\partial M} \;=\; \chi_{\mathrm{mod}}''(d_A^*)\,\frac{\partial d_A^*}{\partial M}.
\]
Solving:
\[
    \frac{\partial d_A^*}{\partial M} \;=\; \frac{\lambda\,\partial^2\Delta\Pi/(\partial d\,\partial M)}{\chi_{\mathrm{mod}}''(d_A^*) - \lambda\,\partial^2\Delta\Pi/\partial d^2}.
\]
At an interior optimum, $\chi_{\mathrm{mod}}'' > 0$ and the denominator is
strictly positive (second-order condition for the maximum). The cross-partial
$\partial^2\Delta\Pi/(\partial d\,\partial M)$ on the increasing branch
is negative (the function $\Delta\Pi$ in $\hat d^{\sigma/2}$ is concave on
its increasing branch, so increases in $M$ reduce the marginal benefit of
further $d$). Hence $\partial d_A^*/\partial M < 0$.
\end{proof}

\subsection{Proof of \autoref{lem:Vv_concave}} \label{app:proof:Vv_concave}

\begin{proof}[Sketch]
Continuity and boundedness follow from standard contraction-mapping
arguments on the Bellman operator over bounded continuous functions on
$[0, \infty)$ (see Stokey, Lucas, and Prescott 1989 for the canonical
treatment).

\textit{Monotonicity of $V^v(\cdot, B)$.} Under Dove with $d_B = 0$, the
within-period payoff is
$U^B(0; M) = \lambda\Pi_B^*(0, S_B; M) + (1-\lambda)u(\bar g_B)$, which is
strictly decreasing in $M$ because $\Pi_B^*$ falls in $M$ (\autoref{lem:contest_M}).
The continuation $M' = (1-\rho)M$ is also strictly increasing in $M$, but
under Dove $V^v(\cdot, B)$ is monotone-decreasing by induction on the
Bellman.

\textit{Single-peakedness of $V^v(\cdot, A)$.} At $M = 0$, $\Pi_A^*(0, S_A; 0) = 1 = \Pi_B^*(0, S_B; 0)$,
so the security premium is zero and the Dove's $u(\bar g_B) > u(\bar g_A)$
implies $V^v(0, A) < V^v(0, B)$ (assuming $\Delta u > 0$, the canonical case).
At intermediate $M$, the security premium activates and the Hawk's $S_A$
advantage delivers $\Pi_A^* > \Pi_B^*$, raising $U^A$ above $U^B$. At very
large $M$, both $\Pi_A^*$ and $\Pi_B^*$ approach 0, and the security
premium $\Delta\Pi$ also vanishes (right tail of the inverted-U); the
Dove's economic advantage $\Delta u$ dominates again. Hence $V^v(\cdot, A)$
is single-peaked.

\textit{Single-crossing wedge.} $\Phi^M(M) = V^v(M, A) - V^v(M, B)$ is
continuous, negative at $M = 0$, positive on an intermediate interval
(if any exists), and crosses zero exactly once on $[0, \infty)$ under
the regularity conditions of \autoref{prop:provocation_M}. The case
$\Phi^M < 0$ everywhere gives $\bar M = \infty$ (no trap).
\end{proof}

\subsection{Proof of \autoref{prop:cutoff}} \label{app:proof:cutoff}

\begin{proof}
By the single-crossing of $\Phi^M$ from \autoref{lem:Vv_concave},
$V^v(M, A) \ge V^v(M, B)$ iff $M \ge \bar M$, hence the voter retains the
Hawk iff $M \ge \bar M$. Symmetrically, retains the Dove iff $V^v(M, B) \ge V^v(M, A)$,
iff $\Phi^M(M) \le 0$, iff $M \le \underline M$. Generic strict inequality
$\underline M < \bar M$ follows from the existence of a hysteresis interval
analogous to Coate--Morris's $[-s, s)$: when the value gap $|\Phi^M|$ is
small, the marginal effect of a regime switch on continuation $M$ depends
on which incumbent is currently in office (Hawk plays $d_A^* > 0$ and
keeps $M$ high; Dove plays $d_B = 0$ and lets $M$ decay). The two thresholds
coincide only in the degenerate case $\rho = 1$ (full instantaneous decay).
\end{proof}

\subsection{Proof of \autoref{prop:MPE_exist}} \label{app:proof:MPE_exist}

\begin{proof}[Sketch]
The Bellman operator $T$ defined by
\[
    (TV)(M, k) \;=\; U^k(d_k(M); M) + \delta\big[\rho^v(M, k)V(M', k) + (1-\rho^v(M, k))V(M', \tilde k)\big]
\]
maps the space $\mathcal{C}_b([0,\infty) \times \{A, B\})$ of bounded
continuous functions to itself, given any measurable strategy profile.
$T$ is a $\delta$-contraction in the sup norm. By Banach's fixed-point
theorem, it has a unique fixed point $V^v(\cdot, \cdot)$.

\textit{Existence of pure-strategy MPE.} Apply Glicksberg's fixed-point
theorem to the joint best-response correspondence on the compact space
of Markov strategies $\{d_k : [0, \bar M^*] \to [0, \bar d^*]\}$, where
$\bar d^*$ is the eventual maximum provocation. Continuity and convexity
of best responses follow from the strict convexity of $\chi_{\mathrm{mod}}$
and concavity of $\Delta\Pi$ on the relevant branch. (Full proof requires
verifying the regularity conditions; we leave details to a longer version.)

\textit{Dove's corner $d_B^\star(M) = 0$.} Under Dove incumbency,
$\chi_{\mathrm{mod}}$ is borne by Dove if $d > 0$. Provoking would move
$\Pi_B^*$ down (\autoref{lem:contest_M}) and thus $V^v(M, B)$ down
(\autoref{lem:Vv_concave}). Hence $d_B^\star = 0$ in any MPE.

\textit{Hawk's $d_A^\star(M)$.} The within-period FOC of \autoref{prop:provocation_M}
(part iv) characterises $d_A^\star$ as a continuous, weakly decreasing
function of $M$. The dynamic best response coincides with the static one
because the Hawk's continuation value $V^A(M', A)$ and the voter's
retention rule are both Markov in $M$, and the Hawk's per-period payoff
is concave in $d$ on the relevant branch.
\end{proof}

\subsection{Proof of \autoref{prop:trap_M}} \label{app:proof:trap_M}

\begin{proof}
Suppose $(M_0, k_0) \in \mathcal{T}$, i.e., $k_0 = A$ and $M_0 \ge \bar M$. We
show $(M_t, k_t) \in \mathcal{T}$ for all $t$.

\textit{Step 1: voter retains.} By $M_0 \ge \bar M$ and \autoref{prop:cutoff},
$\rho^v(M_0, A) = 1$, so $k_1 = A$.

\textit{Step 2: stock stays above cutoff.} Under Hawk with $d_A^\star(M_0) \ge 0$,
the next-period stock is $M_1 = (1-\rho)M_0 + \rho E_m^*(d_A^\star(M_0), S_A; M_0)$.
Two cases:
\begin{itemize}
    \item If $d_A^\star(M_0) > 0$: $E_m^*(\cdot) > 0$, so
    $M_1 \ge (1-\rho)M_0 + \rho \underline E_m$ where
    $\underline E_m := \inf\{E_m^*(d_A^\star(M), S_A; M) : M \ge \bar M\} > 0$.
    Choose $\bar\rho$ small enough that $(1-\bar\rho)\bar M + \bar\rho\underline E_m \ge \bar M$,
    i.e., $\bar\rho \le (\bar M - \bar M)/(\bar M - \underline E_m) = $ a finite positive number
    when $\underline E_m > \bar M$. Otherwise, set
    $\bar\rho \le 1 - \bar M/M_0 + \rho\underline E_m/M_0$, which is positive
    when $M_0$ is sufficiently above $\bar M$.

    \item If $d_A^\star(M_0) = 0$ (paradox-of-restraint regime,
    $M_0 \ge \alpha^2 S_A S_B/\xi$): $E_m^*(0, S_A; M_0) > 0$ because
    $M_0 > 0$, so the contest is still active.
    $M_1 = (1-\rho)M_0 + \rho E_m^*(0, S_A; M_0)$. The stationary level
    of this dynamics satisfies $M^\infty = E_m^*(0, S_A; M^\infty)$, i.e.,
    \[
        M^\infty \;=\; \frac{\beta^2 S_A^{1/2}(\xi M^\infty)^{3/4}}{S_A\alpha + (\xi M^\infty)^{1/2}}.
    \]
    This implicit equation has a unique positive solution
    $M^\infty(S_A, \alpha, \beta, \xi) > 0$. The dynamics converge
    monotonically to $M^\infty$.
\end{itemize}
In either case, $M_1 \ge \bar M$ for $\rho$ small enough, so
$(M_1, A) \in \mathcal{T}$. By induction, $(M_t, A) \in \mathcal{T}$ for all $t$.
\end{proof}

\subsection{Proof of \autoref{prop:cycles_M}} \label{app:proof:cycles_M}

\begin{proof}[Sketch]
Replace $\hat d^\sigma = d^\sigma + \xi M$ with $\hat d^\sigma = d^\sigma + \xi M + \omega$
(absorbing the exogenous threat). Under myopic voting ($\delta \to 0$),
the retention rule reduces to $\rho^v(M, k) = 1$ iff $U^k(d_k^\star; M, \omega) \ge U^{\tilde k}(d_{\tilde k}^\star; M, \omega)$.
Comparative statics in $\omega$, $M$, $\xi$, $\kappa\phi$ follow by
differentiating the Hawk's and Dove's per-period $U$ values and applying
implicit-function theorem to the regime-comparison condition. Recurrent
oscillation under full-support $\omega$ follows from the intermediate-value
theorem on the Markov chain on $(M_t, k_t)$.
\end{proof}

\subsection{Proof of \autoref{cor:paradox_M}} \label{app:proof:paradox_M}

\begin{proof}
Under myopia, retention requires
$U^A(d_A^\star(M); M) + \varepsilon^v_t \ge U^B(0; M)$. A first-order
stochastic decrease in $\tilde v_t = (1 - \Pi_{A,t}^*) + \varepsilon_t^v$
implies a stochastic decrease in $\varepsilon^v_t$ (since $\Pi_{A,t}^*$ is
deterministic given $M$ and $\omega$), which weakly decreases the LHS,
hence weakly increases the probability that the LHS falls below the RHS.
Equivalently, the Hawk's defeat hazard weakly increases.
\end{proof}

\subsection{Proof of \autoref{prop:cs_trap}} \label{app:proof:cs_trap}

\begin{proof}[Sketch]
Differentiate $\Phi^M(\bar M) = 0$ implicitly with respect to each parameter:
\[
    \frac{\partial \bar M}{\partial \theta}
    \;=\; -\,\frac{\partial \Phi^M / \partial \theta}{\partial \Phi^M / \partial M}
    \;\bigg|_{M = \bar M}.
\]
By \autoref{lem:Vv_concave} the denominator $\partial \Phi^M/\partial M > 0$ at $\bar M$.
The numerator's sign depends on the parameter:
\begin{itemize}
    \item $\theta = \rho$: faster decay raises $V^v(M, B)$ (faster recovery
    from inherited grievance under Dove) and lowers $V^v(M, A)$ (Hawk's
    induced $M$ persistence has shorter half-life). Hence
    $\partial \Phi^M/\partial \rho < 0$, giving $\partial \bar M/\partial \rho > 0$.
    \item $\theta = S_A/S_B$: larger Hawk advantage raises $\Delta\Pi$,
    raises $V^v(M, A)$ and lowers $V^v(M, B)$. Hence
    $\partial \Phi^M/\partial(S_A/S_B) > 0$, giving
    $\partial \bar M/\partial(S_A/S_B) < 0$.
    \item $\theta = \lambda$: larger weight on security raises the security
    premium's electoral relevance; analogous to $S_A/S_B$.
    \item $\theta = \Delta u$: larger Dove economic advantage raises
    $V^v(M, B)$, lowers $\Phi^M$, raises $\bar M$.
\end{itemize}
\end{proof}

%==============================================================================
\bibliographystyle{aer}
\bibliography{ref}

\end{document}
